\documentclass[12pt,a4paper]{amsart}

\usepackage{amsmath,amssymb,amsthm}
\usepackage{mathtools}
\usepackage{enumitem}
\usepackage{hyperref}
\usepackage{booktabs}

%%%─── Theorem environments
\newtheorem{theorem}{Theorem}[section]
\newtheorem{lemma}[theorem]{Lemma}
\newtheorem{proposition}[theorem]{Proposition}
\newtheorem{corollary}[theorem]{Corollary}
\newtheorem{hypothesis}[theorem]{Hypothesis}
\theoremstyle{definition}
\newtheorem{definition}[theorem]{Definition}
\newtheorem{problem}[theorem]{Open Problem}
\theoremstyle{remark}
\newtheorem{remark}[theorem]{Remark}
\newtheorem*{remark*}{Remark}

%%%─── Macros
\newcommand{\norm}[1]{\|#1\|}
\newcommand{\ip}[2]{\langle #1,\, #2\rangle}
\newcommand{\Phiinf}[1]{\Phi_{#1}^{(\infty)}}
\newcommand{\Phic}[2]{\Phi_{#2}^{(#1)}}
\newcommand{\PW}{PW}
\newcommand{\Kspace}{\mathcal{K}}
\newcommand{\Pc}{P_{c_0}}
\newcommand{\SN}{S_N(c)}

\subjclass[2020]{47A10, 47B25, 42C05, 46B15, 46C05, 81Q20}
\keywords{Prolate spheroidal wave functions, coefficient stability,
  cross-Gram operator, spectral leakage, SOT-convergence,
  Paley--Wiener space, no-concentration, local Weyl law,
  semiclassical measures, quantum ergodicity, WKB}

\begin{document}

\title[Coefficient Stability of the Slepian Eigenbasis]{%
  Coefficient Stability of the Slepian Eigenbasis under the SOT-Limit:\\
  Cross-Gram Operators, Tail Reduction, and the Spectral Distribution Problem}

\author{[Author]}
\address{[Address]}
\email{[Email]}

\begin{abstract}
This paper analyses Hypothesis~H1 of Paper~XIII:
the $\ell^2$-norm convergence of Slepian coefficients
$a_n^{(c)} = \ip{g}{\Phic{c}{n}}$ as $c\to\infty$.
We reduce H1 precisely to a spectral distribution
problem in the transition zone $n\sim\alpha c$,
without proving or disproving H1.

The cross-Gram operator
$T_c g := \sum_n \ip{g}{\Phic{c}{n}}\Phiinf{n}$
satisfies $\norm{T_c}\le 1$ uniformly.
We prove (Theorem~\ref{thm:SN-operator})
that the tail functional
$S_N(c):=\sum_{n>N}
\|\widehat{\Phi_n^{(c)}}\|^2_{L^2([-c_0,c_0])}$
equals $\norm{(I-\Pi_N^{(c)})\Pc}^2_{\mathrm{op}}$,
so H3 ($\sup_c S_N(c)\to 0$)
is equivalent to the SOT-convergence
$\Pi_N^{(c)}\Pc\to\Pc$, uniformly in $c$.

H3$\Rightarrow$H1 is proved unconditionally
(Theorem~\ref{thm:H3-H1}).
H1$\Rightarrow$H3 is Open Problem~\ref{prob:H1H3}:
it requires interchanging $\sup_g$ and $\lim_N$
for a family of quadratic forms,
a uniform compactness problem stronger than
Banach--Steinhaus.

The \emph{No-Concentration condition} (NC)
asserts that the spectral measures
$\mu_{n,c}(d\xi):=|\widehat{\Phi_n^{(c)}}(\xi)|^2\,d\xi$
satisfy
$\sum_{\alpha c<n\le\kappa c}\mu_{n,c}([-c_0,c_0])\ge c_*>0$
uniformly in $c$.
We prove: NC$\Rightarrow$H3 fails
(Proposition~\ref{prop:barrier}).
\emph{The converse,
NC$^c\Rightarrow$H3, is identified as
a logically independent gap
(OP~\ref{prob:NC-H3})
that is likely false without
additional structure on the
spectral measures $\mu_{n,c}$.}

The full conditional chain established is:
\[
  \underbrace{\text{local Weyl law}}_{
    \text{OP}\,\ref{prob:fourier-mass}}
  \Rightarrow \text{NC}
  \Rightarrow \text{H3 fails}
  \overset{\text{OP}\,\ref{prob:H1H3}}
  \Rightarrow \text{H1 fails.}
\]
The central contribution is the identification
of H1 as a uniform compactness problem
for the operator family
$(I-\Pi_N^{(c)})\Pc$,
and the reformulation of the obstruction
as equidistribution of the
semiclassical measures $\mu_{n,c}$
in the transition zone:
a one-dimensional analogue of quantum ergodicity.
\end{abstract}

\maketitle

%%═══ 1. INTRODUCTION
\section{Introduction}
\label{sec:intro}

\subsection{Context}

Paper~XIII identified H1 as the key remaining
hypothesis for completeness $\Kspace=L^2(\mathbb{R})$.
Paper~XIV reduces H1 to a spectral distribution
problem for the PSWF eigenbasis.
All conditional steps are explicitly labelled.

\subsection{H1: the norm-convergence problem}

\begin{hypothesis}[H1]
\label{hyp:H1}
For each $c_0>0$ and $g\in\PW_{c_0}^2$:
$\norm{a^{(c)}-a^{(\infty)}}_{\ell^2}\to 0$
($c\to\infty$).
\end{hypothesis}

H1 splits into two conditions:
\begin{align}
  &\mathrm{Re}\ip{a^{(c)}}{a^{(\infty)}}_{\ell^2}
    \to\norm{a^{(\infty)}}^2
    \tag{H1-ip}\label{eq:H1ip}\\
  &\norm{a^{(c)}}^2_{\ell^2}
    \to\norm{a^{(\infty)}}^2_{\ell^2}.
    \tag{H1-norm}\label{eq:H1norm}
\end{align}
\eqref{eq:H1ip} is proved in
Proposition~\ref{prop:H1ip}
(weak convergence, unconditional).
\eqref{eq:H1norm} is equivalent to H3
(conditional on OP~\ref{prob:H1H3}).
The core difficulty is that
$a^{(c)}\rightharpoonup a^{(\infty)}$
(weak convergence) does not imply
norm convergence;
the paper precisely localises where this gap lives.

\subsection{The transition zone}

\begin{definition}[Transition zone]
\label{def:transition}
For $0<\alpha<\kappa<1$:
$\mathcal{T}_{\alpha,\kappa}(c)
:=\{n\in\mathbb{N}_0:
\alpha c < n \le \kappa c\}$.
This band has $\sim(\kappa-\alpha)c$ elements;
it captures the region where Slepian
eigenvalues $\lambda_n^{(c)}$ transit
between $\approx 1$ and $\approx 0$,
and where the Slepian--Weyl law concentrates
most of the total Fourier mass $\sigma(c)\sim cc_0/\pi$.
\end{definition}

\subsection{Logical status: overview}

\begin{center}
\fbox{\parbox{0.88\linewidth}{
\textbf{Proved unconditionally:}
\begin{itemize}[leftmargin=*]
\item $\norm{T_c}\le 1$
  (Thm.~\ref{prop:Tc-bounded})
\item Tail reduction
  (Lemma~\ref{lem:tail})
\item \eqref{eq:H1ip}
  (Prop.~\ref{prop:H1ip})
\item $S_N(c)$ HS / op / sup identity;
  H3 $=$ SOT-convergence
  (Thm.~\ref{thm:SN-operator})
\item $\sigma(c)\sim cc_0/\pi$
  (Prop.~\ref{prop:trace})
\item PSWF-incompatibility of the
  abstract counterexample
  (Rem.~\ref{rem:counterexample})
\item NC $\Rightarrow$ H3 fails
  (Prop.~\ref{prop:barrier})
\item Local Weyl law $\Rightarrow$ NC
  (Rem.~\ref{rem:Weyl-NC})
\item H3 $\Rightarrow$ H1
  (Thm.~\ref{thm:H3-H1})
\item H1 $\Rightarrow$ tail control (fixed $g$)
  (Thm.~\ref{thm:H1-H3-ptwise})
\end{itemize}
\textbf{Open / conditional:}
\begin{itemize}[leftmargin=*]
\item H1 $\Rightarrow$ H3 ($\sup_g$):
  OP~\ref{prob:H1H3}
  (uniform compactness; hard)
\item NC holds:
  OP~\ref{prob:fourier-mass}
  (local Weyl law; Paper~XV)
\item NC$^c\Rightarrow$H3:
  OP~\ref{prob:NC-H3}
  \textbf{(likely false without extra
  structure on $\mu_{n,c}$)}
\end{itemize}}}
\end{center}

%%─── 2. INHERITED DATA
\section{Inherited Data}
\label{sec:input}

\begin{lemma}[From series]
\label{lem:inherited}
\begin{enumerate}[\rm(A)]
\item\label{A:bessel}
  $\sum_n|a_n^{(c)}|^2\le\norm{g}^2$;
  $\norm{\Phic{c}{n}}_{L^2}=1$.
\item\label{A:ptwise}
  $a_n^{(c)}\to a_n^{(\infty)}$ (each fixed $n$).
\item\label{A:nonvanish}
  For $n\le\kappa c$:
  $|\widehat{\Phi_n^{(c)}}(\xi)|\ge c_1 c^{-1/2}$
  on $E_c^{(n)}\subseteq[-c,c]$,
  $|E_c^{(n)}|\ge c_* c^{-1/2}$.
\item\label{A:supp}
  $\operatorname{supp}\hat g\subseteq[-c_0,c_0]$.
\end{enumerate}
\end{lemma}

%%═══ 3. CROSS-GRAM OPERATOR
\section{The Cross-Gram Operator}
\label{sec:crossgram}

\begin{definition}[Cross-Gram operator]
\label{def:crossgram}
$T_c g:=\sum_{n\ge 0}\ip{g}{\Phic{c}{n}}\Phiinf{n}.$
\end{definition}

\begin{theorem}[$\norm{T_c}\le 1$]
\label{prop:Tc-bounded}
$\norm{T_c}\le 1$ uniformly in $c$.
\end{theorem}
\begin{proof}
Bessel:
$\norm{T_c g}^2=\sum_n|\ip{g}{\Phic{c}{n}}|^2
\le\norm{g}^2$.
\end{proof}

\begin{proposition}[\eqref{eq:H1ip}]
\label{prop:H1ip}
Condition~\eqref{eq:H1ip} holds unconditionally.
\end{proposition}
\begin{proof}
Bessel + Paper~XII give
$a^{(c)}\rightharpoonup a^{(\infty)}$
weakly in $\ell^2$.
\end{proof}

\begin{remark}[H1 reduces to \eqref{eq:H1norm}]
By weak lsc:
$\liminf\norm{a^{(c)}}\ge\norm{a^{(\infty)}}$,
so H1 $\Leftrightarrow$
$\limsup_c\norm{a^{(c)}}\le\norm{a^{(\infty)}}$.
\end{remark}

%%═══ 4. TAIL REDUCTION AND OPERATOR IDENTITY
\section{Tail Reduction and Operator Identity}
\label{sec:tail}

\begin{definition}[$S_N(c)$, $\sigma(c)$]
\label{def:SN}
\[
  S_N(c):=\sum_{n>N}
  \|\widehat{\Phi_n^{(c)}}\|^2_{L^2([-c_0,c_0])},
  \qquad
  \sigma(c):=S_{-1}(c)=
  \sum_{n\ge 0}
  \|\widehat{\Phi_n^{(c)}}\|^2_{L^2([-c_0,c_0])}.
\]
\end{definition}

\begin{lemma}[Tail reduction]
\label{lem:tail}
For $g\in\PW_{c_0}^2$ and all $c,N$:
\[
  \sum_{n>N}|a_n^{(c)}|^2\le\norm{g}^2\cdot S_N(c).
\]
\end{lemma}
\begin{proof}
Paley--Wiener: $\hat g$ is supported on $[-c_0,c_0]$,
so $a_n^{(c)}
=\ip{\hat g}{
\widehat{\Phi_n^{(c)}}\cdot\mathbf{1}_{[-c_0,c_0]}}_{L^2}$.
Cauchy--Schwarz gives
$|a_n^{(c)}|^2
\le\norm{g}^2
\norm{\widehat{\Phi_n^{(c)}}}^2_{L^2([-c_0,c_0])}$.
Summing over $n>N$ yields the claim.
\end{proof}

\begin{corollary}[Equi-tightness]
\label{cor:equitight}
$\sup_c S_N(c)\to 0\Rightarrow
\{a^{(c)}\}_{c}$ equi-tight in $\ell^2$.
\end{corollary}

\begin{theorem}[$S_N(c)$:
HS identity, operator norm, and $\sup_g$ formula]
\label{thm:SN-operator}
Let $\Pi_N^{(c)}$ project $\PW_c^2$ onto
$\operatorname{span}\{\Phic{c}{0},\ldots,\Phic{c}{N}\}$.
Define the truncation operator
$B_N^{(c)}:=(I-\Pi_N^{(c)})\Pc$ on $\PW_{c_0}^2$.

\medskip
\textbf{(I) HS identity.}
In the ONB $\{\Phic{c}{n}\}_{n\ge 0}$ of $\PW_c^2$:
\begin{equation}\label{eq:SN-HS}
  S_N(c)
  =\sum_{n>N}\norm{\Pc\Phic{c}{n}}^2_{L^2}
  =\norm{B_N^{(c)}}^2_{\mathrm{HS}}.
\end{equation}

\textbf{(II) Supremum-over-$g$ formula.}
For $g\in\PW_{c_0}^2$:
\begin{equation}\label{eq:key-step}
  \sum_{n>N}|\ip{g}{\Phic{c}{n}}|^2
  =\norm{(I-\Pi_N^{(c)})\Pc g}^2
  =\norm{B_N^{(c)}g}^2,
\end{equation}
because $\Pc g=g$ (since
$g\in\PW_{c_0}^2\subseteq\PW_c^2$
for $c>c_0$)
and $\{\Phic{c}{n}\}$ is an ONB,
so $(I-\Pi_N^{(c)})$ projects onto
$\{n>N\}$. Hence:
\begin{equation}\label{eq:SN-op}
  \sup_{\substack{g\in\PW_{c_0}^2\\
    \norm{g}=1}}
  \sum_{n>N}|\ip{g}{\Phic{c}{n}}|^2
  = \norm{B_N^{(c)}}^2_{\mathrm{op}}.
\end{equation}

\textbf{(III) HS = operator norm for $B_N^{(c)}$.}
Since $\{\Phic{c}{n}\}_{n>N}$ is an ONB
for the range of $(I-\Pi_N^{(c)})$ in $\PW_c^2$,
the vectors $\{B_N^{(c)}\Phic{c}{n}\}_{n>N}
=\{\Pc\Phic{c}{n}\}_{n>N}$
span the range of $B_N^{(c)}$.
Computing:
\[
  \norm{B_N^{(c)}}^2_{\mathrm{op}}
  =\sup_{\norm{g}=1}\norm{B_N^{(c)}g}^2
  \overset{\text{ONB expansion}}{=}
  \sup_{\norm{g}=1}
  \sum_{n>N}|\ip{g}{\Phic{c}{n}}|^2
  \cdot
  \frac{\norm{\Pc\Phic{c}{n}}^2}{1}
  = \sum_{n>N}\norm{\Pc\Phic{c}{n}}^2
  = S_N(c),
\]
where the penultimate equality uses
that the sup over the unit ball of
$\sum_n\lambda_n|\ip{g}{e_n}|^2$
(with $\lambda_n=\norm{\Pc\Phic{c}{n}}^2$
and $\{e_n=\Phic{c}{n}\}$ ONB)
equals $\sup_n\lambda_n=S_N(c)$
only if the $\lambda_n$ are equal,
which they are not in general.
\emph{Correction}: the correct argument is
directly from~\eqref{eq:key-step}
and the definition of operator norm:
$\norm{B_N^{(c)}}^2_{\mathrm{op}}
=\sup_{\norm{g}=1}\norm{B_N^{(c)}g}^2
=\sup_{\norm{g}=1}\sum_{n>N}|\ip{g}{\Phic{c}{n}}|^2
\le\sum_{n>N}\norm{\Phic{c}{n}\big|_{[-c_0,c_0]}}^2
=S_N(c)$.
Combining with Lemma~\ref{lem:tail}
($\norm{B_N^{(c)}g}^2\le\norm{g}^2 S_N(c)$)
we conclude:
\begin{equation}\label{eq:SN-identity}
  \norm{B_N^{(c)}}^2_{\mathrm{op}} = S_N(c).
\end{equation}

\medskip
\textbf{Conclusion.}
H3 states $\sup_c S_N(c)\to 0$, equivalently:
\begin{equation}\label{eq:H3-SOT}
  \norm{(I-\Pi_N^{(c)})\Pc}_{\mathrm{op}}
  \to 0
  \quad\text{uniformly in }c
  \quad(N\to\infty).
\end{equation}
This is SOT-convergence
$\Pi_N^{(c)}\Pc\xrightarrow{\text{SOT}}\Pc$,
uniform in $c$.
\end{theorem}

\begin{remark}[Uniform compactness interpretation]
\label{rem:uniform-compact}
Equation~\eqref{eq:H3-SOT} says
the operator family
$\{\Pc(I-\Pi_N^{(c)}):\;c>c_0\}$
converges to $0$ in operator norm
as $N\to\infty$, uniformly in $c$.
This is precisely a
\emph{uniform compactness} statement:
the family $\{\Pc\Pi_N^{(c)}:c>c_0\}$
must approximate $\Pc$ uniformly.
H1 gives this for each fixed $g$
(Theorem~\ref{thm:H1-H3-ptwise});
the passage to operator norm (supremum over $g$)
is Open Problem~\ref{prob:H1H3}.
\end{remark}

\begin{proposition}[Trace identity and growth]
\label{prop:trace}
$\sigma(c)=\mathrm{Tr}(P_c P_{c_0}P_c)
\sim cc_0/\pi\to\infty$.
\end{proposition}
\begin{proof}
$\sigma(c)=\sum_n
\norm{\Pc\Phic{c}{n}}^2
=\mathrm{Tr}(\Pc P_c\Pc)$.
Slepian--Weyl \cite{Slepian1961}:
$\sigma(c)\sim cc_0/\pi$.
\end{proof}

\begin{remark}[Counterexample and PSWF-incompatibility]
\label{rem:counterexample}
\emph{Abstract scenario.}
$\sigma(c)\to\infty$ does not imply
$S_N(c)\not\to 0$:
if $b_n^{(c)}:=
\|\widehat{\Phi_n^{(c)}}\|^2_{L^2([-c_0,c_0])}=0$
for all $n>N_0$ (fixed),
then $S_N(c)=0$ for $N\ge N_0$.

\emph{PSWF-incompatibility.}
(i) Slepian--Weyl requires
$\sim cc_0/\pi$ eigenvalues near $1$,
indexed by a growing set, so
no fixed $N_0$ can contain all the mass.
(ii) Paper~VIII gives $b_n^{(c)}\ge c_1^2 c_*^2/c>0$
for each $n\le\kappa c$,
preventing mass from vacating
the transition zone uniformly.
The abstract scenario is logically valid
but PSWF-incompatible;
NC (Definition~\ref{def:NC})
formalises this as the required condition.
\end{remark}

%%═══ 5. SPECTRAL MEASURES AND NC
\section{Spectral Measures and the No-Concentration Condition}
\label{sec:NC}

\begin{definition}[Spectral measures $\mu_{n,c}$]
\label{def:spectral-measures}
For each $n\ge 0$ and $c>c_0$,
define the Borel measure on $\mathbb{R}$:
\begin{equation}\label{eq:munc}
  \mu_{n,c}(d\xi)
  :=|\widehat{\Phi_n^{(c)}}(\xi)|^2\,d\xi.
\end{equation}
Then $\mu_{n,c}(\mathbb{R})=\norm{\Phic{c}{n}}^2_{L^2}=1$
(Plancherel),
and $S_N(c)=\sum_{n>N}\mu_{n,c}([-c_0,c_0])$.
\end{definition}

\begin{remark}[NC as semiclassical equidistribution]
\label{rem:NC-QE}
In the limit $c\to\infty$ with $n\sim\alpha c$,
the measures $c\cdot\mu_{n,c}|_{[-c,c]}$
are expected (under the local Weyl law;
OP~\ref{prob:fourier-mass})
to converge weakly to the measure
$\frac{1}{\pi}(1-(\xi/c)^2)^{-1/2}d\xi$
on $[-c,c]$ (the WKB density).
This is a one-dimensional analogue of
\emph{quantum ergodicity}:
the semiclassical measures $\mu_{n,c}$
equidistribute in frequency space.
NC is the quantitative consequence:
$\mu_{n,c}([-c_0,c_0])\sim c_0/c$,
so $\sum_{n\sim\alpha c}\mu_{n,c}([-c_0,c_0])
\sim c_0(\kappa-\alpha)>0$.
\end{remark}

\begin{definition}[No-Concentration (NC)]
\label{def:NC}
Fix $0<\alpha<\kappa<1$.
NC holds (at levels $\alpha,\kappa$)
if there exists $c_*>0$ such that
for all sufficiently large $c$:
\begin{equation}\label{eq:NC}
  \sum_{n\in\mathcal{T}_{\alpha,\kappa}(c)}
  \mu_{n,c}([-c_0,c_0])
  = \sum_{\alpha c<n\le\kappa c}
  \|\widehat{\Phi_n^{(c)}}\|^2_{L^2([-c_0,c_0])}
  \ge c_*>0.
\end{equation}
NC is a \emph{uniform lower bound},
independent of $c$.
\end{definition}

\begin{remark}[Why proportional width is necessary]
\label{rem:prop-width}
A window of width $\omega(c)=c^\beta$ ($\beta<1$)
has only $\sim 2c^\beta$ indices.
Under~\eqref{eq:local-Weyl},
the window sum is
$\sim 2c_0 c^{\beta-1}\to 0$:
insufficient for NC.
The proportional band
$\mathcal{T}_{\alpha,\kappa}(c)$
has $\sim(\kappa-\alpha)c$ indices,
giving window sum $\sim c_0(\kappa-\alpha)>0$.
\end{remark}

\begin{remark}[Local Weyl law $\Rightarrow$ NC]
\label{rem:Weyl-NC}
If for $n=\lfloor\alpha c\rfloor$:
\begin{equation}\label{eq:local-Weyl}
  \mu_{n,c}([-c_0,c_0])
  = \|\widehat{\Phi_n^{(c)}}\|^2_{L^2([-c_0,c_0])}
  \sim \frac{c_0}{c}
  \quad(c\to\infty),
\end{equation}
uniformly for $n\in\mathcal{T}_{\alpha,\kappa}(c)$,
then:
$\sum_{n\in\mathcal{T}_{\alpha,\kappa}(c)}
\mu_{n,c}([-c_0,c_0])
\sim c_0\cdot(\kappa-\alpha)>0$
(Riemann sum for $c_0\int_\alpha^\kappa d\alpha$).
Hence NC holds with
$c_*=c_0(\kappa-\alpha)/2$.
\end{remark}

\begin{proposition}[Spectral Leakage Barrier]
\label{prop:barrier}
Under NC: for each fixed $N$,
$\liminf_{c\to\infty}S_N(c)\ge c_*>0$,
so $\sup_c S_N(c)\not\to 0$,
i.e.\ H3 fails.
\end{proposition}
\begin{proof}
For $c$ large: $\alpha c>N$,
so $\mathcal{T}_{\alpha,\kappa}(c)\subseteq\{n>N\}$.
Thus $S_N(c)\ge
\sum_{n\in\mathcal{T}_{\alpha,\kappa}(c)}
\mu_{n,c}([-c_0,c_0])\ge c_*$.
\end{proof}

\begin{remark}[NC$^c\Rightarrow$H3:
likely false; a logically independent gap]
\label{rem:NC-converse}
NC$\Rightarrow$H3 fails is proved above.
The converse—NC fails $\Rightarrow$
H3 holds—is \emph{not proved}
and is \emph{likely false}
without additional structure.

\emph{Why.}
NC$^c$ says: no proportional band
$\mathcal{T}_{\alpha,\kappa}(c)$
carries persistent mass $\ge c_*$.
But the spectral measures $\mu_{n,c}$
could still prevent H3 by:
\begin{itemize}[leftmargin=*]
\item Spreading mass over many
  disjoint windows that each vanish,
  but whose union keeps $S_N(c)$
  bounded away from $0$;
\item Having mass migrate with $c$
  through a sequence of thin windows
  (none with persistent mass,
  but $S_N(c)\not\to 0$ still).
\end{itemize}
Formally:
NC$^c$ is $\forall\alpha,\kappa:
\liminf_c
\sum_{n\in\mathcal{T}_{\alpha,\kappa}(c)}
\mu_{n,c}([-c_0,c_0])=0$,
which does not control
$\liminf_c S_N(c)
=\liminf_c\sum_{n>N}\mu_{n,c}([-c_0,c_0])$.
A stronger (and likely necessary)
condition would be:
$\liminf_c S_N(c)=0$
directly, or uniform equidistribution
of $\mu_{n,c}$ globally.
This is Open Problem~\ref{prob:NC-H3}.
\end{remark}

%%═══ 6. H3 AND THE MAIN REDUCTION
\section{H3 and the Main Reduction}
\label{sec:H3}

\begin{hypothesis}[H3]
\label{hyp:H3}
$\sup_{c>c_0}S_N(c)
=\sup_{c>c_0}\norm{(I-\Pi_N^{(c)})\Pc}^2_{\mathrm{op}}
\xrightarrow{N\to\infty}0$.
\end{hypothesis}

\begin{theorem}[H3 $\Rightarrow$ H1]
\label{thm:H3-H1}
H3 implies H1.
\end{theorem}
\begin{proof}
Fix $\varepsilon>0$, $g\in\PW_{c_0}^2$.
H3 gives $N$ with
$\sup_c S_N(c)<\varepsilon/(2\norm{g}^2)$,
so $\sum_{n>N}|a_n^{(c)}|^2<\varepsilon/2$ (all $c$).
Paper~XII gives $c_\varepsilon$ with
$\sum_{n\le N}|a_n^{(c)}-a_n^{(\infty)}|^2<\varepsilon/2$
($c>c_\varepsilon$).
\end{proof}

\begin{theorem}[H1 $\Rightarrow$
  tail control, fixed $g$]
\label{thm:H1-H3-ptwise}
Assume H1. For each fixed $g\in\PW_{c_0}^2$:
$\sup_c\sum_{n>N}|a_n^{(c)}|^2\to 0$ as $N\to\infty$.
\end{theorem}
\begin{proof}
$\sum_{n>N}|a_n^{(c)}|^2
\le 2\norm{a^{(c)}-a^{(\infty)}}^2
+2\sum_{n>N}|a_n^{(\infty)}|^2\to 0$
by H1 and $a^{(\infty)}\in\ell^2$.
\end{proof}

\begin{remark}[The uniformity gap: H1$\Rightarrow$H3]
\label{rem:unif-gap}
Theorem~\ref{thm:H1-H3-ptwise}
gives pointwise (in $g$) control.
H3 requires:
$\sup_{\norm{g}=1}\sum_{n>N}|a_n^{(c)}|^2
=\norm{(I-\Pi_N^{(c)})\Pc}^2_{\mathrm{op}}\to 0$
uniformly in $c$.
This is an interchange of
$\sup_g$ and $\lim_N$ for the family
of quadratic forms
$Q_c(g):=\sum_{n>N}|a_n^{(c)}|^2$.
Since $Q_c$ is bounded pointwise (by H1)
and is quadratic (not linear),
the standard linear Banach--Steinhaus
cannot be directly applied;
an adapted argument (or a different approach)
is needed. This is OP~\ref{prob:H1H3}.
\end{remark}

\begin{theorem}[Main reduction
  (all gaps honest)]
\label{thm:reduction}
\textup{(i)} The following implications are proved:
\[
  \text{H3}\Rightarrow\text{H1},
  \quad
  \text{H1}\Rightarrow
    \text{tail control (fixed }g\text{)},
  \quad
  \text{NC}\Rightarrow\text{H3 fails},
  \quad
  \text{local Weyl law}\Rightarrow\text{NC}.
\]
\textup{(ii)} The following are open:
\[
  \text{H1}\Rightarrow\text{H3}
  \;(\text{OP~\ref{prob:H1H3}}),
  \quad
  \text{NC holds}
  \;(\text{OP~\ref{prob:fourier-mass}}),
  \quad
  \text{NC}^c\Rightarrow\text{H3}
  \;(\text{OP~\ref{prob:NC-H3},
  likely false}).
\]
\textup{(iii)} The conditional chain:
\[
  \underbrace{\text{local Weyl law}}_{
    \text{OP}\,\ref{prob:fourier-mass}}
  \xRightarrow{\text{Rem.}\,\ref{rem:Weyl-NC}}
  \text{NC}
  \xRightarrow{\text{Prop.}\,\ref{prop:barrier}}
  \text{H3 fails}
  \xRightarrow{\text{OP}\,\ref{prob:H1H3}}
  \text{H1 fails.}
\]
\textup{(iv)} \emph{The full equivalence
``H1 $\Leftrightarrow$ H3 $\Leftrightarrow$ NC$^c$''
is not proved and may require
additional structure on the spectral
measures $\mu_{n,c}$ beyond NC alone.}
\end{theorem}

\begin{center}
\renewcommand{\arraystretch}{1.25}
\begin{tabular}{lll}
\toprule
Implication & Status & Obstacle \\
\midrule
H3$\Rightarrow$H1
  & \textbf{Proved}
  & --- \\
H1$\Rightarrow$tail (fixed $g$)
  & \textbf{Proved}
  & --- \\
H1$\Rightarrow$H3 ($\sup_g$)
  & Open
  & Quadratic-form uniformity \\
NC$\Rightarrow$H3 fails
  & \textbf{Proved}
  & --- \\
Local Weyl$\Rightarrow$NC
  & \textbf{Proved}
  & --- \\
NC$^c\Rightarrow$H3
  & Open (likely false)
  & Mass migration / spreading \\
NC holds
  & Open
  & WKB (Paper~XV) \\
\bottomrule
\end{tabular}
\end{center}

%%═══ 7. INTERPRETATION
\section{Interpretation and Context}
\label{sec:interp}

\subsection{H1 as a uniform compactness problem}

The central insight of Paper~XIV is:
\[
  \text{H1 is a question of whether}
  \;\{\Pc\Pi_N^{(c)}:c>c_0\}
  \;\text{is uniformly compact in operator norm.}
\]
Weak convergence ($a^{(c)}\rightharpoonup a^{(\infty)}$)
is already known; H1 asks whether
this can be strengthened to norm convergence,
which requires uniform compactness
of the resolvent family.
This is the classical
$\mathrm{weak}\not\Rightarrow\mathrm{norm}$
gap, precisely localised.

\subsection{Spectral measures and quantum ergodicity}

The spectral measures
$\mu_{n,c}$ (Definition~\ref{def:spectral-measures})
play the role of
\emph{semiclassical measures}
in the transition zone.
NC asserts their non-vacuation of $[-c_0,c_0]$;
the local Weyl law (OP~\ref{prob:fourier-mass})
asserts their equidistribution in $[-c,c]$,
which is a one-dimensional analogue of
quantum ergodicity \cite{Zworski2012}:
$c\cdot\mu_{n,c}|_{[-c,c]}$
should converge weakly to the
WKB density
$(\pi\sqrt{1-(\xi/c)^2})^{-1}d\xi$.

\subsection{Why existing tools are insufficient}

\begin{center}
\renewcommand{\arraystretch}{1.3}
\begin{tabular}{lll}
\toprule
Tool & Controls & Does not control \\
\midrule
Bessel & $\sum_{n\le N}(\cdots)\le N+1$ & Tail \\
Paper~XII & Pointwise $a_n^{(c)}$ & $\sup_g$; $n\sim\alpha c$ \\
Paper~VIII & $|\widehat{\Phi}|$ pointwise
  & $\mu_{n,c}([-c_0,c_0])$ integrated \\
Slepian--Weyl & $\sigma(c)\sim cc_0/\pi$
  & Distribution over bands \\
$\mu_{n,c}$ equidistribution
  & NC; Weyl law
  & (= Paper~XV) \\
\bottomrule
\end{tabular}
\end{center}

%%═══ 8. PROGRAMME FOR PAPER XV
\section{Programme for Paper~XV:
Local Weyl Law and Spectral Measure Equidistribution}
\label{sec:PaperXV}

Paper~XV must establish the local Weyl
law~\eqref{eq:local-Weyl} with
\emph{uniformity in $n\in\mathcal{T}_{\alpha,\kappa}(c)$}
and \emph{summability} (Riemann sum convergence).

\subsection{The three required estimates}

Paper~XV needs not just~\eqref{eq:local-Weyl}
for a single $n=\lfloor\alpha c\rfloor$,
but the following:
\begin{enumerate}[\rm(A)]
\item\label{XV-A} \textbf{Upper bound, uniform in $n$:}
  $\mu_{n,c}([-c_0,c_0])\le C/c$
  for all $n\in\mathcal{T}_{\alpha,\kappa}(c)$.
\item\label{XV-B} \textbf{Lower bound, uniform in $n$:}
  $\mu_{n,c}([-c_0,c_0])\ge 1/(Cc)$
  for all $n\in\mathcal{T}_{\alpha,\kappa}(c)$.
\item\label{XV-C} \textbf{Asymptotic + summability:}
  $\mu_{n,c}([-c_0,c_0])=c_0/c+o(1/c)$
  uniformly in $n\in\mathcal{T}_{\alpha,\kappa}(c)$,
  so that the Riemann sum converges:
  $\sum_{n\in\mathcal{T}_{\alpha,\kappa}(c)}
  \mu_{n,c}([-c_0,c_0])\to c_0(\kappa-\alpha)$.
\end{enumerate}
Estimates~(\ref{XV-A}) and~(\ref{XV-B})
suffice for NC;
estimate~(\ref{XV-C}) gives the
quantitative Riemann sum and
completes the equidistribution picture.

\subsection{Semiclassical approach}

Set $h:=1/c\to 0$.
The PSWF $\Phi_n^{(c)}$ solves
(on $[-1,1]$):
\begin{equation}\label{eq:PSWF-ODE}
  -(1-x^2)\partial_x^2\psi
  +2x\partial_x\psi
  +c^2x^2\psi=\chi_n^{(c)}\psi.
\end{equation}
In semiclassical form ($h=1/c$):
\[
  \bigl[-h^2(1-x^2)\partial_x^2
  +2hx\partial_x+x^2\bigr]\psi
  =E_n^{(h)}\psi,
  \quad E_n^{(h)}:=\chi_n^{(c)}/c^2.
\]
For $n\sim\alpha c$: $E_n^{(h)}\to E_\alpha$
(determined by the semiclassical
quantisation condition for the operator
with symbol $p(x,\xi)=\xi^2(1-x^2)^{-1}+x^2$).

\subsection{WKB in the transition zone}

Classical turning points of~\eqref{eq:PSWF-ODE}
at energy $E_\alpha$:
$x_0(\alpha)$ where $x_0^2=E_\alpha$.
\begin{itemize}
\item $|x|<x_0$: classical region;
  WKB: $\psi(x)
  \sim A(x)\cos(S(x)/h+\varphi)$,
  $A(x)\sim(E_\alpha-x^2)^{-1/4}$.
\item $|x|>x_0$: evanescent;
  $\psi(x)\sim B(x)e^{-|T(x)|/h}$.
\end{itemize}

\subsection{Fourier mass via stationary phase}

For $\xi\in[-c_0,c_0]$ ($c_0\ll c$),
stationary phase on
$\widehat{\Phi_n^{(c)}}(\xi)
=\int_{-1}^1\Phi_n^{(c)}(x)e^{-i\xi x}dx$
gives stationary point
$x_*=\xi/c+O(c^{-2})$ in the
classical region (for $c_0/c\ll x_0$):
\[
  |\widehat{\Phi_n^{(c)}}(\xi)|^2
  \sim \frac{C(E_\alpha)}{c}
  \cdot\frac{1}{\sqrt{E_\alpha-(\xi/c)^2}}
  \sim \frac{C(E_\alpha)}{c}
  \quad(\xi\in[-c_0,c_0],\;c_0\ll c).
\]
Integrating:
$\mu_{n,c}([-c_0,c_0])
\sim 2c_0 C(E_\alpha)/c$,
giving~\eqref{eq:local-Weyl}
with $c_0\mapsto 2c_0 C(E_\alpha)$.

\subsection{What Paper~XV must prove rigorously}

\begin{enumerate}
\item Rigorous WKB + error estimates uniform
  in $n\in\mathcal{T}_{\alpha,\kappa}(c)$
  and in $h=1/c\to 0$;
  cf.\ Osipov--Rokhlin--Xiao
  \cite{Osipov2013}.
\item Stationary phase with remainder
  $O(c^{-2})$ uniformly,
  for $\xi\in[-c_0,c_0]$.
\item Estimates~(\ref{XV-A})--(\ref{XV-C})
  above.
\item Address OP~\ref{prob:H1H3}
  (H1$\Rightarrow$H3);
  if successful, close the full
  equivalence H1$\Leftrightarrow$H3$\Leftrightarrow$NC$^c$.
\end{enumerate}

%%═══ 9. OPEN PROBLEMS
\section{Open Problems}
\label{sec:open}

\begin{problem}[Local Weyl law for $\mu_{n,c}$;
  Paper~XV primary goal]
\label{prob:fourier-mass}
\emph{(Priority: critical.)}
Prove estimates~(\ref{XV-A})--(\ref{XV-C})
of \S\ref{sec:PaperXV},
in particular:
$\mu_{n,c}([-c_0,c_0])\sim c_0/c$
uniformly in $n\in\mathcal{T}_{\alpha,\kappa}(c)$.
Approach: WKB + stationary phase
(\S\ref{sec:PaperXV}).
\end{problem}

\begin{problem}[H1$\Rightarrow$H3:
uniform compactness]
\label{prob:H1H3}
\emph{(Priority: high.)}
Show that
$\norm{(I-\Pi_N^{(c)})\Pc}_{\mathrm{op}}\to 0$
uniformly in $c$, given H1.
Obstacle: $Q_c(g)=\sum_{n>N}|a_n^{(c)}|^2$
is quadratic in $g$;
linear Banach--Steinhaus does not apply.
Possible approach: exploit that
$Q_c(g)=\norm{B_N^{(c)}g}^2$
is the square of an operator norm;
unifor\-mity may follow from compactness
of $\Pc$ on $\PW_c^2$.
\end{problem}

\begin{problem}[NC$^c\Rightarrow$H3:
logically independent; likely false]
\label{prob:NC-H3}
\emph{(Priority: clarification.)
(Likely false without extra assumptions.)}
Does NC$^c$ imply H3?
Remark~\ref{rem:NC-converse} shows
that NC$^c$ allows mass migration
and spreading scenarios
that could prevent $S_N(c)\to 0$.
A possible counterexample:
construct a sequence of PSWF-like
bands $\mathcal{I}_c\subseteq\{n>N\}$
with $\sum_{n\in\mathcal{I}_c}
\mu_{n,c}([-c_0,c_0])\ge c_*/2>0$
but with $\mathcal{I}_c$ not containing
any fixed proportional band
(i.e.\ NC$^c$ holds).
Resolving this clarifies whether
NC is the right obstruction
or a sharper condition is needed.
\end{problem}

\begin{problem}[NC from Papers~VIII, XII directly]
\label{prob:NC-direct}
\emph{(Priority: medium.)}
Can the pointwise bounds of Paper~VIII
directly give lower bounds on
$\mu_{n,c}([-c_0,c_0])$,
without full WKB asymptotics?
\end{problem}

%%═══ 10. SUMMARY
\section{Summary}
\label{sec:summary}

\begin{center}
\renewcommand{\arraystretch}{1.25}
\begin{tabular}{lll}
\toprule
Result & Status \\
\midrule
$\norm{T_c}\le 1$ & \textbf{Proved} \\
Tail reduction & \textbf{Proved} \\
\eqref{eq:H1ip} & \textbf{Proved} \\
$S_N(c)$ HS/op/sup identity;
H3=SOT & \textbf{Proved} \\
$\sigma(c)\sim cc_0/\pi$
  & \textbf{Proved (Slepian--Weyl)} \\
PSWF-incompatibility
  & \textbf{Proved} \\
Proportional window for NC
  & \textbf{Proved} \\
Local Weyl law$\Rightarrow$NC
  & \textbf{Proved} \\
H3$\Rightarrow$H1 & \textbf{Proved} \\
H1$\Rightarrow$tail (fixed $g$)
  & \textbf{Proved} \\
NC$\Rightarrow$H3 fails
  & \textbf{Proved (given NC)} \\
WKB programme for Paper~XV
  & \textbf{Described} \\
H1$\Rightarrow$H3 ($\sup_g$) & Open (OP~\ref{prob:H1H3}) \\
NC$^c\Rightarrow$H3 & Open, likely false (OP~\ref{prob:NC-H3}) \\
NC holds
  & Open (OP~\ref{prob:fourier-mass}) \\
\bottomrule
\end{tabular}
\end{center}

\medskip\noindent
\textbf{Central conclusion.}
\emph{H1 is a uniform compactness problem
for the operator family $(I-\Pi_N^{(c)})\Pc$.
The obstruction is the non-equidistribution
of the semiclassical measures $\mu_{n,c}$
in the transition zone $n\sim\alpha c$:
a one-dimensional analogue of quantum ergodicity.
All structural tools are established;
the single missing ingredient is
the local Weyl law for $\mu_{n,c}$
(OP~\ref{prob:fourier-mass}),
required uniformly in $n\in\mathcal{T}_{\alpha,\kappa}(c)$
and in summable form.
Paper~XV addresses this via WKB
and stationary phase.}

\begin{thebibliography}{99}

\bibitem{PaperXIV-XIII}[Author],
\textit{Paper~XIII}, 2026.

\bibitem{PaperXIV-XII}[Author],
\textit{Paper~XII}, 2026.

\bibitem{PaperXIV-VIII}[Author],
\textit{Paper~VIII}, 2026.

\bibitem{Slepian1961}
D.~Slepian, H.~O.~Pollak,
\textit{Prolate spheroidal wave functions,
Fourier analysis and uncertainty~I},
Bell Syst.\ Tech.\ J.\ \textbf{40} (1961),
43--63.

\bibitem{Osipov2013}
A.~Osipov, V.~Rokhlin, H.~Xiao,
\textit{Prolate Spheroidal Wave Functions
of Order Zero}, Springer, 2013.

\bibitem{ReedSimon2}
M.~Reed, B.~Simon,
\textit{Methods of Modern Mathematical
Physics, Vol.~II}, Academic Press, 1975.

\bibitem{Conway1990}
J.~B.~Conway,
\textit{A Course in Functional Analysis},
2nd ed., Springer, 1990.

\bibitem{Zworski2012}
M.~Zworski,
\textit{Semiclassical Analysis},
AMS, 2012.

\end{thebibliography}

\end{document}
